\section{Source concordance}
\label{supp:concordance}

\Cref{tab:concordance} lists every numbered item of the main paper with its entry, its role, the
line of the theorem catalog of the development at which that entry begins, and its Lean declarations
with their kind, file and line. Declaration names omit the common prefix of the library; names
beginning with \path{Instances.} belong to the instance modules. Items argued on paper only are
marked as such.

{\scriptsize
\begin{longtable}{@{}l>{\raggedright\arraybackslash}p{2.8cm}>{\raggedright\arraybackslash}p{10.2cm}@{}}
\caption{Source concordance.}\label{tab:concordance}\\
\toprule
\textbf{Item} & \textbf{Entry and role} & \textbf{Lean declarations} \\
\midrule
\endfirsthead
\toprule
\textbf{Item} & \textbf{Entry and role} & \textbf{Lean declarations} \\
\midrule
\endhead
\bottomrule
\endfoot
\Cref{def:theory-package} & setting, definition & \path{TheoryPackage} (structure) \newline \path{lean/SixBirdsFoundationsV/Definitional/ESystem.lean}:16 \\
\Cref{def:d1} & D1, definition \newline catalog line 54 & \path{Refines} (def) \newline \path{lean/SixBirdsFoundationsV/Definitional/RepairJoin.lean}:13 \newline \path{FiberEquiv} (def) \newline \path{lean/SixBirdsFoundationsV/Definitional/RepairJoin.lean}:16 \newline \path{repairJoin} (def) \newline \path{lean/SixBirdsFoundationsV/Definitional/RepairJoin.lean}:19 \\
\Cref{prop:d1} & D1, general theorem \newline catalog line 54 & \path{join_refines_left} (theorem) \newline \path{lean/SixBirdsFoundationsV/Definitional/RepairJoin.lean}:23 \newline \path{join_refines_right} (theorem) \newline \path{lean/SixBirdsFoundationsV/Definitional/RepairJoin.lean}:29 \newline \path{vee_greatest_lower_bound} (theorem) \newline \path{lean/SixBirdsFoundationsV/Definitional/RepairJoin.lean}:35 \newline \path{join_idem} (theorem) \newline \path{lean/SixBirdsFoundationsV/Definitional/RepairJoin.lean}:57 \newline \path{join_comm} (theorem) \newline \path{lean/SixBirdsFoundationsV/Definitional/RepairJoin.lean}:66 \newline \path{join_assoc} (theorem) \newline \path{lean/SixBirdsFoundationsV/Definitional/RepairJoin.lean}:75 \newline \path{join_well_defined} (theorem) \newline \path{lean/SixBirdsFoundationsV/Definitional/RepairJoin.lean}:42 \\
\Cref{def:d2} & D2, definition \newline catalog line 56 & \path{predictiveSurplusValue} (def) \newline \path{lean/SixBirdsFoundationsV/Definitional/PredictiveSurplus.lean}:55 \newline \path{predictiveSurplusXi} (def) \newline \path{lean/SixBirdsFoundationsV/Definitional/PredictiveSurplus.lean}:91 \\
\Cref{prop:d2} & D2, conditional general theorem \newline catalog line 56 & \path{predictiveSurplusValue_nonnegative_iff} (theorem) \newline \path{lean/SixBirdsFoundationsV/Definitional/PredictiveSurplus.lean}:58 \newline \path{predictiveSurplusValue_identical} (theorem) \newline \path{lean/SixBirdsFoundationsV/Definitional/PredictiveSurplus.lean}:72 \newline \path{predictiveSurplusValue_exact_baseline} (theorem) \newline \path{lean/SixBirdsFoundationsV/Definitional/PredictiveSurplus.lean}:78 \newline \path{predictiveSurplusXi_nonnegative_of_psd} (theorem) \newline \path{lean/SixBirdsFoundationsV/Definitional/PredictiveSurplus.lean}:96 \newline \path{predictiveSurplusXi_chain_rule_monotone} (theorem) \newline \path{lean/SixBirdsFoundationsV/Definitional/PredictiveSurplus.lean}:103 \newline \path{predictiveSurplusXi_same_family_saturation} (theorem) \newline \path{lean/SixBirdsFoundationsV/Definitional/PredictiveSurplus.lean}:136 \\
\Cref{def:d3} & D3, definition \newline catalog line 60 & \path{FineSourceTag} (inductive) \newline \path{lean/SixBirdsFoundationsV/Definitional/CarriedRecord.lean}:15 \newline \path{CarriedSource} (def) \newline \path{lean/SixBirdsFoundationsV/Definitional/CarriedRecord.lean}:26 \newline \path{CarriedRecord} (def) \newline \path{lean/SixBirdsFoundationsV/Definitional/CarriedRecord.lean}:54 \newline \path{CarriedInstrument} (def) \newline \path{lean/SixBirdsFoundationsV/Definitional/CarriedRecord.lean}:89 \\
\Cref{lem:d3} & D3, consistency lemma \newline catalog line 60 & \path{not_carriedSource_fallback} (theorem) \newline \path{lean/SixBirdsFoundationsV/Definitional/CarriedRecord.lean}:115 \newline \path{not_carriedSource_independent_pair_witness} (theorem) \newline \path{lean/SixBirdsFoundationsV/Definitional/CarriedRecord.lean}:122 \newline \path{not_carriedSource_simulation_trace} (theorem) \newline \path{lean/SixBirdsFoundationsV/Definitional/CarriedRecord.lean}:131 \newline \path{not_carriedSource_ablation_record} (theorem) \newline \path{lean/SixBirdsFoundationsV/Definitional/CarriedRecord.lean}:139 \newline \path{not_carriedSource_unknown} (theorem) \newline \path{lean/SixBirdsFoundationsV/Definitional/CarriedRecord.lean}:147 \newline \path{not_carriedSource_contradictory} (theorem) \newline \path{lean/SixBirdsFoundationsV/Definitional/CarriedRecord.lean}:154 \newline \path{not_carriedRecord_of_step_not_in_scope} (theorem) \newline \path{lean/SixBirdsFoundationsV/Definitional/CarriedRecord.lean}:198 \newline \path{not_carriedInstrument_of_incomplete_inventory} (theorem) \newline \path{lean/SixBirdsFoundationsV/Definitional/CarriedRecord.lean}:265 \\
\Cref{def:d4} & D4, definition \newline catalog line 62 & \path{RepairSort} (def) \newline \path{lean/SixBirdsFoundationsV/Definitional/ESystem.lean}:13 \newline \path{CarriedLedger} (structure) \newline \path{lean/SixBirdsFoundationsV/Definitional/ESystem.lean}:72 \newline \path{RepairMove} (structure) \newline \path{lean/SixBirdsFoundationsV/Definitional/ESystem.lean}:85 \newline \path{ActiveCarriedInstrument} (structure) \newline \path{lean/SixBirdsFoundationsV/Definitional/ESystem.lean}:97 \newline \path{ESystem} (structure) \newline \path{lean/SixBirdsFoundationsV/Definitional/ESystem.lean}:177 \\
\Cref{def:d4-step} & D4, definition \newline catalog line 62 & \path{LawfulRepairStep} (structure) \newline \path{lean/SixBirdsFoundationsV/Definitional/ESystem.lean}:123 \newline \path{ESystem.RepairStep} (def) \newline \path{lean/SixBirdsFoundationsV/Definitional/ESystem.lean}:198 \\
\Cref{lem:d4} & D4, consistency lemma \newline catalog line 62 & \path{lawfulRepairStep_requires_detection} (theorem) \newline \path{lean/SixBirdsFoundationsV/Definitional/ESystem.lean}:319 \newline \path{lawfulRepairStep_requires_gate} (theorem) \newline \path{lean/SixBirdsFoundationsV/Definitional/ESystem.lean}:346 \newline \path{lawfulRepairStep_requires_kernel_transition} (theorem) \newline \path{lean/SixBirdsFoundationsV/Definitional/ESystem.lean}:373 \newline \path{lawfulRepairStep_requires_reaudit} (theorem) \newline \path{lean/SixBirdsFoundationsV/Definitional/ESystem.lean}:400 \newline \path{not_lawfulRepairStep_without_detection} (theorem) \newline \path{lean/SixBirdsFoundationsV/Definitional/ESystem.lean}:427 \newline \path{not_lawfulRepairStep_without_gate} (theorem) \newline \path{lean/SixBirdsFoundationsV/Definitional/ESystem.lean}:453 \newline \path{not_lawfulRepairStep_without_kernel_transition} (theorem) \newline \path{lean/SixBirdsFoundationsV/Definitional/ESystem.lean}:479 \newline \path{not_lawfulRepairStep_without_reaudit} (theorem) \newline \path{lean/SixBirdsFoundationsV/Definitional/ESystem.lean}:505 \newline \path{no_repairSort_p6} (theorem) \newline \path{lean/SixBirdsFoundationsV/Definitional/ESystem.lean}:215 \\
\Cref{def:d5} & D5, definition \newline catalog line 66 & \path{ChallengeEpisode} (structure) \newline \path{lean/SixBirdsFoundationsV/Definitional/ClosedLoopScope.lean}:77 \newline \path{ChallengeHistory} (structure) \newline \path{lean/SixBirdsFoundationsV/Definitional/ClosedLoopScope.lean}:88 \newline \path{ClosedLoopScope} (def) \newline \path{lean/SixBirdsFoundationsV/Definitional/ClosedLoopScope.lean}:109 \\
\Cref{lem:d5} & D5, consistency lemma \newline catalog line 66 & \path{closedLoopScope_complete} (theorem) \newline \path{lean/SixBirdsFoundationsV/Definitional/ClosedLoopScope.lean}:128 \newline \path{closedLoopScope_entry_carried} (theorem) \newline \path{lean/SixBirdsFoundationsV/Definitional/ClosedLoopScope.lean}:146 \newline \path{not_closedLoopScope_of_listed_source_fallback} (theorem) \newline \path{lean/SixBirdsFoundationsV/Definitional/ClosedLoopScope.lean}:474 \newline \path{not_closedLoopScope_of_listed_source_unknown} (theorem) \newline \path{lean/SixBirdsFoundationsV/Definitional/ClosedLoopScope.lean}:496 \newline \path{not_closedLoopScope_of_listed_source_contradictory} (theorem) \newline \path{lean/SixBirdsFoundationsV/Definitional/ClosedLoopScope.lean}:518 \newline \path{not_closedLoopScope_of_listed_generatedByS_false} (theorem) \newline \path{lean/SixBirdsFoundationsV/Definitional/ClosedLoopScope.lean}:540 \newline \path{not_closedLoopScope_of_listed_inScope_false} (theorem) \newline \path{lean/SixBirdsFoundationsV/Definitional/ClosedLoopScope.lean}:562 \newline \path{fallback_entry_cannot_borrow_carriedness} (theorem) \newline \path{lean/SixBirdsFoundationsV/Definitional/ClosedLoopScope.lean}:584 \\
\Cref{def:d6} & D6, definition \newline catalog line 68 & \path{ProbeEconomy} (structure) \newline \path{lean/SixBirdsFoundationsV/Definitional/ProbeEconomy.lean}:45 \newline \path{SameFamilySaturated} (def) \newline \path{lean/SixBirdsFoundationsV/Definitional/ProbeEconomy.lean}:30 \newline \path{AcquisitionStrict} (def) \newline \path{lean/SixBirdsFoundationsV/Definitional/ProbeEconomy.lean}:34 \newline \path{LawfulAllocation} (structure) \newline \path{lean/SixBirdsFoundationsV/Definitional/ProbeEconomy.lean}:117 \newline \path{LawfulAcquisition} (structure) \newline \path{lean/SixBirdsFoundationsV/Definitional/ProbeEconomy.lean}:157 \newline \path{LawfulRetirement} (structure) \newline \path{lean/SixBirdsFoundationsV/Definitional/ProbeEconomy.lean}:198 \\
\Cref{lem:d6} & D6, consistency lemma \newline catalog line 68 & \path{lawfulAllocation_no_new_active_probe} (theorem) \newline \path{lean/SixBirdsFoundationsV/Definitional/ProbeEconomy.lean}:290 \newline \path{lawfulAcquisition_requires_acquisitionStrict} (theorem) \newline \path{lean/SixBirdsFoundationsV/Definitional/ProbeEconomy.lean}:228 \newline \path{lawfulRetirement_requires_active} (theorem) \newline \path{lean/SixBirdsFoundationsV/Definitional/ProbeEconomy.lean}:313 \newline \path{lawfulRetirement_probe_absent_after} (theorem) \newline \path{lean/SixBirdsFoundationsV/Definitional/ProbeEconomy.lean}:333 \newline \path{not_lawfulAcquisition_of_sameFamilySaturated} (theorem) \newline \path{lean/SixBirdsFoundationsV/Definitional/ProbeEconomy.lean}:249 \newline \path{xi_sameFamilySaturated_not_acquisitionStrict} (theorem) \newline \path{lean/SixBirdsFoundationsV/Definitional/ProbeEconomy.lean}:395 \\
\Cref{def:device} & E3, E2, instance data \newline catalog line 133 & \path{Instances.eventTheory} (def) \newline \path{lean/SixBirdsFoundationsV/Instances/E3Executed.lean}:6 \newline \path{Instances.eventPolicy} (def) \newline \path{lean/SixBirdsFoundationsV/Instances/E3Executed.lean}:34 \newline \path{Instances.eventInstrument} (def) \newline \path{lean/SixBirdsFoundationsV/Instances/E3Executed.lean}:48 \newline \path{Instances.eventSystem} (def) \newline \path{lean/SixBirdsFoundationsV/Instances/E3Executed.lean}:61 \\
\Cref{prop:device-step} & D4, instance \newline catalog line 62 & \path{Instances.eventAt} (theorem) \newline \path{lean/SixBirdsFoundationsV/Instances/E3Executed.lean}:38 \newline \path{Instances.eventCarried} (theorem) \newline \path{lean/SixBirdsFoundationsV/Instances/E3Executed.lean}:44 \newline \path{Instances.eventRepair} (theorem) \newline \path{lean/SixBirdsFoundationsV/Instances/E3Executed.lean}:81 \\
\Cref{def:episode-linked} & E3, definition \newline catalog line 133 & \path{MaintenanceReinstatementFor} (def) \newline \path{lean/SixBirdsFoundationsV/Laws/E3SelfMaintainingReclosure.lean}:133 \newline \path{Instances.EpisodeLinkedReinstatement} (def) \newline \path{lean/SixBirdsFoundationsV/Instances/E3Episode.lean}:5 \newline \path{Instances.FullyHistoryLinkedReinstatement} (def) \newline \path{lean/SixBirdsFoundationsV/Instances/E3Episode.lean}:79 \\
\Cref{thm:e3-episode} & E3, instance \newline catalog line 133 & \path{Instances.eventTrajectory} (theorem) \newline \path{lean/SixBirdsFoundationsV/Instances/E3Executed.lean}:20 \newline \path{Instances.eventTransition} (theorem) \newline \path{lean/SixBirdsFoundationsV/Instances/E3Executed.lean}:29 \newline \path{Instances.eventHistory} (def) \newline \path{lean/SixBirdsFoundationsV/Instances/E3Executed.lean}:97 \newline \path{Instances.eventApp} (def) \newline \path{lean/SixBirdsFoundationsV/Instances/E3Executed.lean}:106 \newline \path{Instances.eventMaintenance} (def) \newline \path{lean/SixBirdsFoundationsV/Instances/E3Executed.lean}:116 \newline \path{Instances.eventReinstatement} (def) \newline \path{lean/SixBirdsFoundationsV/Instances/E3Executed.lean}:122 \newline \path{Instances.eventE3Reinstatement} (theorem) \newline \path{lean/SixBirdsFoundationsV/Instances/E3Executed.lean}:151 \newline \path{Instances.eventHistoryAgrees} (theorem) \newline \path{lean/SixBirdsFoundationsV/Instances/E3Executed.lean}:162 \newline \path{Instances.eventFullyHistoryLinked} (theorem) \newline \path{lean/SixBirdsFoundationsV/Instances/E3Executed.lean}:170 \newline \path{Instances.eventExecutedReinstatement} (theorem) \newline \path{lean/SixBirdsFoundationsV/Instances/E3Executed.lean}:187 \newline \path{Instances.noEpisodeLinkedWithEmptyHistory} (theorem) \newline \path{lean/SixBirdsFoundationsV/Instances/E3Episode.lean}:62 \\
\Cref{def:declared-measure} & E2, definition \newline catalog line 109 & \path{CarriedInstrumentLevel} (structure) \newline \path{lean/SixBirdsFoundationsV/Laws/E2BoundedReflexivity.lean}:20 \newline \path{CarriedInstrumentLevelOccurrence} (def) \newline \path{lean/SixBirdsFoundationsV/Laws/E2BoundedReflexivity.lean}:40 \newline \path{Instances.DeclaredCapacityMeasure} (structure) \newline \path{lean/SixBirdsFoundationsV/Instances/E2Declared.lean}:6 \newline \path{Instances.DeclaredOccurrence} (def) \newline \path{lean/SixBirdsFoundationsV/Instances/E2Declared.lean}:43 \newline \path{Instances.DeclaredFootprint} (def) \newline \path{lean/SixBirdsFoundationsV/Instances/E2Declared.lean}:54 \\
\Cref{thm:e2-declared} & E2, general theorem \newline catalog line 109 & \path{Instances.declaredRecordList_nodup} (theorem) \newline \path{lean/SixBirdsFoundationsV/Instances/E2Declared.lean}:84 \newline \path{Instances.declaredRecordList_subset} (theorem) \newline \path{lean/SixBirdsFoundationsV/Instances/E2Declared.lean}:104 \newline \path{Instances.declaredRecordList_length} (theorem) \newline \path{lean/SixBirdsFoundationsV/Instances/E2Declared.lean}:116 \newline \path{Instances.declaredCapacityBound} (theorem) \newline \path{lean/SixBirdsFoundationsV/Instances/E2Declared.lean}:131 \newline \path{Instances.declaredTowerCapacityBound} (theorem) \newline \path{lean/SixBirdsFoundationsV/Instances/E2Declared.lean}:148 \newline \path{Instances.copiedRecordReindexNotDeclared} (theorem) \newline \path{lean/SixBirdsFoundationsV/Instances/E2Declared.lean}:168 \\
\Cref{thm:e2-tower} & E2, instance \newline catalog line 109 & \path{Instances.eventCheckedSystem} (def) \newline \path{lean/SixBirdsFoundationsV/Instances/E2EventTower.lean}:19 \newline \path{Instances.eventCheckedSystemRepair} (theorem) \newline \path{lean/SixBirdsFoundationsV/Instances/E2EventTower.lean}:39 \newline \path{Instances.eventInspectLower} (def) \newline \path{lean/SixBirdsFoundationsV/Instances/E2EventTower.lean}:67 \newline \path{Instances.eventLowerInspectionPasses} (theorem) \newline \path{lean/SixBirdsFoundationsV/Instances/E2EventTower.lean}:70 \newline \path{Instances.eventUpperRecordedCheck} (theorem) \newline \path{lean/SixBirdsFoundationsV/Instances/E2EventTower.lean}:96 \newline \path{Instances.eventCheckedMeasure} (def) \newline \path{lean/SixBirdsFoundationsV/Instances/E2EventTower.lean}:128 \newline \path{Instances.eventCheckedTowerAuditsLowerStack} (theorem) \newline \path{lean/SixBirdsFoundationsV/Instances/E2EventTower.lean}:234 \newline \path{Instances.eventCheckedLevelsDisjoint} (theorem) \newline \path{lean/SixBirdsFoundationsV/Instances/E2EventTower.lean}:248 \newline \path{Instances.eventCheckedConcreteCapacity} (theorem) \newline \path{lean/SixBirdsFoundationsV/Instances/E2EventTower.lean}:259 \newline \path{Instances.eventCheckedCapacityBound} (theorem) \newline \path{lean/SixBirdsFoundationsV/Instances/E2EventTower.lean}:263 \newline \path{Instances.eventCheckedStatusSaturated} (theorem) \newline \path{lean/SixBirdsFoundationsV/Instances/E2EventTower.lean}:282 \newline \path{Instances.eventCheckedExecutedRepair} (theorem) \newline \path{lean/SixBirdsFoundationsV/Instances/E2EventTower.lean}:288 \\
\Cref{prop:e2-self} & E2, imported, weaker form \newline catalog line 109 & \path{E2_SelfSoundnessObstruction} (theorem) \newline \path{lean/SixBirdsFoundationsV/Laws/E2BoundedReflexivity.lean}:831 \\
\Cref{prop:e2-v1-vacuous} & E2, revision note: version-1 measure vacuous \newline catalog line 109 & \path{AuditTowerCapacityMeasure} (structure) \newline \path{lean/SixBirdsFoundationsV/Laws/E2BoundedReflexivity.lean}:267 \newline \path{E2_CapacityBound} (theorem) \newline \path{lean/SixBirdsFoundationsV/Laws/E2BoundedReflexivity.lean}:1325 \newline \path{Instances.noE2WithOccurringLevel} (theorem) \newline \path{lean/SixBirdsFoundationsV/Instances/SpikeModel.lean}:104 \newline \path{Instances.spikeNoE2} (theorem) \newline \path{lean/SixBirdsFoundationsV/Instances/SpikeModel.lean}:89 \\
\Cref{thm:e6-kkt} & E6, general theorem \newline catalog line 205 & \path{KKTWitness} (structure) \newline \path{lean/SixBirdsFoundationsV/Laws/E6E9PricedAccess.lean}:109 \newline \path{GenuineScarcity} (def) \newline \path{lean/SixBirdsFoundationsV/Laws/E6E9PricedAccess.lean}:41 \newline \path{E6_AttentionKKT} (theorem) \newline \path{lean/SixBirdsFoundationsV/Laws/E6E9PricedAccess.lean}:265 \newline \path{E6_SlackCollapse} (theorem) \newline \path{lean/SixBirdsFoundationsV/Laws/E6E9PricedAccess.lean}:372 \\
\Cref{thm:e6-caps} & E6, general theorem \newline catalog line 205 & \path{Instances.weightedCappedObjective} (def) \newline \path{lean/SixBirdsFoundationsV/Instances/E6AtCaps.lean}:6 \newline \path{Instances.cappedFeasible} (def) \newline \path{lean/SixBirdsFoundationsV/Instances/E6Capped.lean}:9 \newline \path{Instances.saturatedSlope} (def) \newline \path{lean/SixBirdsFoundationsV/Instances/E6AtCaps.lean}:9 \newline \path{Instances.weightedSaturatedSupportingPlane} (theorem) \newline \path{lean/SixBirdsFoundationsV/Instances/E6AtCaps.lean}:13 \newline \path{Instances.CapMultiplierKKT} (structure) \newline \path{lean/SixBirdsFoundationsV/Instances/E6AtCaps.lean}:45 \newline \path{Instances.capMultiplierKKT_sufficient} (theorem) \newline \path{lean/SixBirdsFoundationsV/Instances/E6AtCaps.lean}:69 \\
\Cref{thm:e6-atcap} & E6, numerical instance \newline catalog line 205 & \path{Instances.capExampleKKT} (def) \newline \path{lean/SixBirdsFoundationsV/Instances/E6AtCaps.lean}:109 \newline \path{Instances.capExampleSlopeZero} (theorem) \newline \path{lean/SixBirdsFoundationsV/Instances/E6AtCaps.lean}:128 \newline \path{Instances.capExampleGloballyOptimal} (theorem) \newline \path{lean/SixBirdsFoundationsV/Instances/E6AtCaps.lean}:137 \newline \path{Instances.unitCapExampleGloballyOptimal} (theorem) \newline \path{lean/SixBirdsFoundationsV/Instances/E6AtCaps.lean}:145 \\
\Cref{def:e1-forcing} & E1, definition \newline catalog line 80 & \path{Instances.ActualFamilyForcing} (structure) \newline \path{lean/SixBirdsFoundationsV/Instances/E1Periodic.lean}:8 \newline \path{GenericallyNovelChallenge} (structure) \newline \path{lean/SixBirdsFoundationsV/Laws/E1Internalization.lean}:190 \newline \path{StrictSelfExtension} (structure) \newline \path{lean/SixBirdsFoundationsV/Laws/E1Internalization.lean}:198 \newline \path{InfinitelyManyStrictSelfExtensions} (def) \newline \path{lean/SixBirdsFoundationsV/Laws/E1Internalization.lean}:202 \\
\Cref{prop:e1-forcing} & E1, general theorem \newline catalog line 80 & \path{Instances.actualFamilyStrictExtensions} (theorem) \newline \path{lean/SixBirdsFoundationsV/Instances/E1Periodic.lean}:17 \\
\Cref{thm:e1-clocked} & E1, instance \newline catalog line 80 & \path{Instances.clockedTheory} (def) \newline \path{lean/SixBirdsFoundationsV/Instances/E1Periodic.lean}:125 \newline \path{Instances.clockedSystem} (def) \newline \path{lean/SixBirdsFoundationsV/Instances/E1Periodic.lean}:183 \newline \path{Instances.clockedInstalls} (def) \newline \path{lean/SixBirdsFoundationsV/Instances/E1Periodic.lean}:204 \newline \path{Instances.clockedOddRepair} (theorem) \newline \path{lean/SixBirdsFoundationsV/Instances/E1Periodic.lean}:210 \newline \path{Instances.clockedEvenPerturbation} (theorem) \newline \path{lean/SixBirdsFoundationsV/Instances/E1Periodic.lean}:222 \newline \path{Instances.clockedFamilyEveryEntryInstalled} (theorem) \newline \path{lean/SixBirdsFoundationsV/Instances/E1Periodic.lean}:343 \newline \path{Instances.clockedUnbounded} (theorem) \newline \path{lean/SixBirdsFoundationsV/Instances/E1Periodic.lean}:363 \newline \path{Instances.periodicNovel} (theorem) \newline \path{lean/SixBirdsFoundationsV/Instances/E1Periodic.lean}:286 \newline \path{Instances.periodicStrictAt} (theorem) \newline \path{lean/SixBirdsFoundationsV/Instances/E1Periodic.lean}:290 \newline \path{Instances.clockedForcing} (def) \newline \path{lean/SixBirdsFoundationsV/Instances/E1Periodic.lean}:371 \newline \path{Instances.clockedStrictExtensions} (theorem) \newline \path{lean/SixBirdsFoundationsV/Instances/E1Periodic.lean}:380 \newline \path{Instances.periodicRefinementsDistinct} (theorem) \newline \path{lean/SixBirdsFoundationsV/Instances/E1Periodic.lean}:296 \\
\Cref{prop:e1-v1-vacuous} & E1, revision note: version-1 certificate uninhabited \newline catalog line 80 & \path{FiniteForcingStrictnessCertified} (structure) \newline \path{lean/SixBirdsFoundationsV/Laws/E1Internalization.lean}:235 \newline \path{E1_StrictSelfExtension} (theorem) \newline \path{lean/SixBirdsFoundationsV/Laws/E1Internalization.lean}:1753 \newline \path{Instances.finiteForcingNoPerpetualNovelty} (theorem) \newline \path{lean/SixBirdsFoundationsV/Instances/E1ForcingVacuity.lean}:7 \newline \path{Instances.finiteForcingUninhabited} (theorem) \newline \path{lean/SixBirdsFoundationsV/Instances/E1ForcingVacuity.lean}:40 \newline \path{Instances.E1_StrictSelfExtension_noCertificate} (theorem) \newline \path{lean/SixBirdsFoundationsV/Instances/E1ForcingVacuity.lean}:64 \\
\Cref{def:e12-setting} & E12, definition \newline catalog line 343 & \path{Instances.E12Semantic.PartClosure} (structure) \newline \path{lean/SixBirdsFoundationsV/Instances/E12Semantic.lean}:13 \newline \path{Instances.E12Semantic.System} (structure) \newline \path{lean/SixBirdsFoundationsV/Instances/E12Semantic.lean}:20 \newline \path{Instances.E12Semantic.RepairClosed} (def) \newline \path{lean/SixBirdsFoundationsV/Instances/E12Semantic.lean}:30 \newline \path{Instances.E12Semantic.SelfMaintaining} (def) \newline \path{lean/SixBirdsFoundationsV/Instances/E12Semantic.lean}:33 \newline \path{Instances.E12Semantic.Candidate} (def) \newline \path{lean/SixBirdsFoundationsV/Instances/E12Semantic.lean}:37 \newline \path{Instances.E12Semantic.Maximal} (def) \newline \path{lean/SixBirdsFoundationsV/Instances/E12Semantic.lean}:40 \newline \path{Instances.E12Semantic.Integrated} (def) \newline \path{lean/SixBirdsFoundationsV/Instances/E12Semantic.lean}:43 \\
\Cref{prop:e12-structure} & E12, general theorem \newline catalog line 343 & \path{Instances.E12Semantic.repairClosed_inter} (theorem) \newline \path{lean/SixBirdsFoundationsV/Instances/E12Semantic.lean}:47 \newline \path{Instances.E12Semantic.integrated_disjoint_or_equal} (theorem) \newline \path{lean/SixBirdsFoundationsV/Instances/E12Semantic.lean}:67 \\
\Cref{thm:e12-repair} & E12, instance \newline catalog line 343 & \path{Instances.E12Semantic.genuineRepairClosure} (def) \newline \path{lean/SixBirdsFoundationsV/Instances/E12Semantic.lean}:219 \newline \path{Instances.E12Semantic.repairSystem} (def) \newline \path{lean/SixBirdsFoundationsV/Instances/E12Semantic.lean}:257 \newline \path{Instances.E12Semantic.repairPart_integrated} (theorem) \newline \path{lean/SixBirdsFoundationsV/Instances/E12Semantic.lean}:288 \newline \path{Instances.E12Semantic.repairedIndividualWitness} (theorem) \newline \path{lean/SixBirdsFoundationsV/Instances/E12Semantic.lean}:342 \newline \path{Instances.E12Semantic.healthyOnly_closure_loadBearing} (theorem) \newline \path{lean/SixBirdsFoundationsV/Instances/E12Semantic.lean}:360 \\
\Cref{prop:e12-overlap} & E12, negative instance \newline catalog line 343 & \path{Instances.E12Semantic.secondSystem} (def) \newline \path{lean/SixBirdsFoundationsV/Instances/E12Semantic.lean}:145 \newline \path{Instances.E12Semantic.secondLeft_maximal} (theorem) \newline \path{lean/SixBirdsFoundationsV/Instances/E12Semantic.lean}:161 \newline \path{Instances.E12Semantic.secondRight_maximal} (theorem) \newline \path{lean/SixBirdsFoundationsV/Instances/E12Semantic.lean}:174 \newline \path{Instances.E12Semantic.second_no_integrated} (theorem) \newline \path{lean/SixBirdsFoundationsV/Instances/E12Semantic.lean}:187 \\
\Cref{def:e13-transport} & E13, definition \newline catalog line 367 & \path{Instances.E13Semantic.System} (structure) \newline \path{lean/SixBirdsFoundationsV/Instances/E13Semantic.lean}:6 \newline \path{Instances.E13Semantic.Interface} (structure) \newline \path{lean/SixBirdsFoundationsV/Instances/E13Semantic.lean}:14 \newline \path{Instances.E13Semantic.RepairTransport} (structure) \newline \path{lean/SixBirdsFoundationsV/Instances/E13Semantic.lean}:24 \\
\Cref{thm:e13-chain} & E13, general theorem \newline catalog line 367 & \path{Instances.E13Semantic.compose} (def) \newline \path{lean/SixBirdsFoundationsV/Instances/E13Semantic.lean}:33 \newline \path{Instances.E13Semantic.compose_payment} (theorem) \newline \path{lean/SixBirdsFoundationsV/Instances/E13Semantic.lean}:49 \newline \path{Instances.E13Semantic.RepairChain} (inductive) \newline \path{lean/SixBirdsFoundationsV/Instances/E13Semantic.lean}:151 \newline \path{Instances.E13Semantic.chain_obstruction_bound} (theorem) \newline \path{lean/SixBirdsFoundationsV/Instances/E13Semantic.lean}:159 \newline \path{Instances.E13Semantic.chain_length_bound} (theorem) \newline \path{lean/SixBirdsFoundationsV/Instances/E13Semantic.lean}:170 \\
\Cref{thm:e13-cross} & E13, instance \newline catalog line 367 & \path{Instances.E13Semantic.crossTransport} (def) \newline \path{lean/SixBirdsFoundationsV/Instances/E13Semantic.lean}:185 \newline \path{Instances.E13Semantic.cross_context_zero_to_one} (theorem) \newline \path{lean/SixBirdsFoundationsV/Instances/E13Semantic.lean}:203 \newline \path{Instances.E13Semantic.cross_context_one_to_zero} (theorem) \newline \path{lean/SixBirdsFoundationsV/Instances/E13Semantic.lean}:206 \newline \path{Instances.E13Semantic.cross_context_two_step_bound} (theorem) \newline \path{lean/SixBirdsFoundationsV/Instances/E13Semantic.lean}:216 \newline \path{Instances.E13Semantic.sourceToReceiver} (def) \newline \path{lean/SixBirdsFoundationsV/Instances/E13Semantic.lean}:73 \newline \path{Instances.E13Semantic.receiverToThird} (def) \newline \path{lean/SixBirdsFoundationsV/Instances/E13Semantic.lean}:88 \newline \path{Instances.E13Semantic.concrete_composite_payment} (theorem) \newline \path{lean/SixBirdsFoundationsV/Instances/E13Semantic.lean}:111 \newline \path{Instances.E13Semantic.concrete_composite_reduces} (theorem) \newline \path{lean/SixBirdsFoundationsV/Instances/E13Semantic.lean}:115 \\
\Cref{prop:e13-flat} & E13, negative instance \newline catalog line 367 & \path{Instances.E13Semantic.roleOnlyInterface} (def) \newline \path{lean/SixBirdsFoundationsV/Instances/E13Semantic.lean}:126 \newline \path{Instances.E13Semantic.roleOnly_mapsRepair} (theorem) \newline \path{lean/SixBirdsFoundationsV/Instances/E13Semantic.lean}:134 \newline \path{Instances.E13Semantic.roleOnly_no_reduction} (theorem) \newline \path{lean/SixBirdsFoundationsV/Instances/E13Semantic.lean}:139 \\
\Cref{ex:e6-nonconcave} & E6, counterexample \newline catalog line 205 & argued on paper; not formalized \\
\Cref{rem:e6-unlinked} & E6, counterexample \newline catalog line 205 & argued on paper; not formalized \\
\Cref{prop:e6-sufficient} & E6, sufficiency, argued on paper \newline catalog line 205 & argued on paper; not formalized \\
\Cref{thm:e9} & E9, general theorem \newline catalog line 273 & \path{E9_CuriosityArgmax} (theorem) \newline \path{lean/SixBirdsFoundationsV/Laws/E6E9PricedAccess.lean}:479 \newline \path{E6_E9_AccessArbitration} (theorem) \newline \path{lean/SixBirdsFoundationsV/Laws/E6E9PricedAccess.lean}:587 \newline \path{E6_E9_AccessArbitration_epsilon} (theorem) \newline \path{lean/SixBirdsFoundationsV/Laws/E6E9PricedAccess.lean}:620 \newline \path{selected_access_move_ratio_bound} (theorem) \newline \path{lean/SixBirdsFoundationsV/Laws/E6E9PricedAccess.lean}:661 \newline \path{selected_access_move_ratio_bound_epsilon} (theorem) \newline \path{lean/SixBirdsFoundationsV/Laws/E6E9PricedAccess.lean}:695 \\
\Cref{prop:e9-saturation} & E9, general theorem \newline catalog line 273 & \path{XiAcqDischarge} (def) \newline \path{lean/SixBirdsFoundationsV/Laws/E6E9PricedAccess.lean}:406 \newline \path{E9_SameFamilyStrictness} (theorem) \newline \path{lean/SixBirdsFoundationsV/Laws/E6E9PricedAccess.lean}:414 \\
\Cref{prop:e15-identity} & E15, general theorem \newline catalog line 412 & \path{DerivedOfflineBudgetGeometry} (structure) \newline \path{lean/SixBirdsFoundationsV/Laws/E15OfflineReclosure.lean}:691 \newline \path{kappaOn} (def) \newline \path{lean/SixBirdsFoundationsV/Laws/E15OfflineReclosure.lean}:722 \newline \path{kappaOff} (def) \newline \path{lean/SixBirdsFoundationsV/Laws/E15OfflineReclosure.lean}:732 \newline \path{freedAllocation} (def) \newline \path{lean/SixBirdsFoundationsV/Laws/E15OfflineReclosure.lean}:727 \newline \path{Instances.offlineCapacityFromExhaustion} (theorem) \newline \path{lean/SixBirdsFoundationsV/Instances/E15BudgetIdentity.lean}:6 \newline \path{Instances.geometryCapacityFromExhaustion} (theorem) \newline \path{lean/SixBirdsFoundationsV/Instances/E15BudgetIdentity.lean}:29 \\
\Cref{prop:e4-brittle} & E4, general theorem \newline catalog line 157 & \path{E4_Brittleness} (theorem) \newline \path{lean/SixBirdsFoundationsV/Laws/E4RepairCompilation.lean}:816 \\
\Cref{tmpl:e1} & E1, conditional classification \newline catalog line 80 & \path{E1_StatusPartition} (theorem) \newline \path{lean/SixBirdsFoundationsV/Laws/E1Internalization.lean}:1164 \newline \path{E1_Internalization} (theorem) \newline \path{lean/SixBirdsFoundationsV/Laws/E1Internalization.lean}:1525 \newline \path{E1_NoGeneratorDichotomy} (theorem) \newline \path{lean/SixBirdsFoundationsV/Laws/E1Internalization.lean}:1770 \\
\Cref{tmpl:e3} & E3, conditional classification \newline catalog line 133 & \path{E3_TwoLevelFixedPoint} (theorem) \newline \path{lean/SixBirdsFoundationsV/Laws/E3SelfMaintainingReclosure.lean}:2275 \newline \path{E3_StatusPartition} (theorem) \newline \path{lean/SixBirdsFoundationsV/Laws/E3SelfMaintainingReclosure.lean}:2418 \newline \path{E3_RegressStoppedByE2} (theorem) \newline \path{lean/SixBirdsFoundationsV/Laws/E3SelfMaintainingReclosure.lean}:2795 \newline \path{E3_AblationSeparation} (theorem) \newline \path{lean/SixBirdsFoundationsV/Laws/E3SelfMaintainingReclosure.lean}:2963 \\
\Cref{tmpl:e4} & E4, certificate specification \newline catalog line 157 & \path{compilationLawful_requiredCoreGatesPass} (theorem) \newline \path{lean/SixBirdsFoundationsV/Laws/E4RepairCompilation.lean}:190 \newline \path{E4_Compilation} (theorem) \newline \path{lean/SixBirdsFoundationsV/Laws/E4RepairCompilation.lean}:774 \newline \path{E4_StatusPartition} (theorem) \newline \path{lean/SixBirdsFoundationsV/Laws/E4RepairCompilation.lean}:843 \\
\Cref{def:e5} & E5, definition \newline catalog line 181 & \path{CompleteCollapseRescueInventory} (structure) \newline \path{lean/SixBirdsFoundationsV/Laws/E5ReclosureCollapse.lean}:756 \newline \path{CollapsedCore} (def) \newline \path{lean/SixBirdsFoundationsV/Laws/E5ReclosureCollapse.lean}:920 \\
\Cref{lem:e5} & E5, consistency lemma \newline catalog line 181 & \path{E5_CollapseFalsifier} (theorem) \newline \path{lean/SixBirdsFoundationsV/Laws/E5ReclosureCollapse.lean}:2034 \newline \path{E5_ReclosureCollapse} (theorem) \newline \path{lean/SixBirdsFoundationsV/Laws/E5ReclosureCollapse.lean}:2041 \newline \path{E5_RevivedExcludesLowerPriority} (theorem) \newline \path{lean/SixBirdsFoundationsV/Laws/E5ReclosureCollapse.lean}:2618 \newline \path{E5_IrreversibleCollapse} (theorem) \newline \path{lean/SixBirdsFoundationsV/Laws/E5ReclosureCollapse.lean}:2720 \newline \path{E5_SubsidyWithdrawalReclassification} (theorem) \newline \path{lean/SixBirdsFoundationsV/Laws/E5ReclosureCollapse.lean}:2794 \\
\Cref{tmpl:e7} & E7, conditional classification \newline catalog line 229 & \path{E7_AlarmTrichotomy} (theorem) \newline \path{lean/SixBirdsFoundationsV/Laws/E7Alarm.lean}:653 \newline \path{E7_PreemptionSignature} (theorem) \newline \path{lean/SixBirdsFoundationsV/Laws/E7Alarm.lean}:949 \\
\Cref{def:e8} & E8, definition \newline catalog line 253 & \path{ControlPriceComponent} (structure) \newline \path{lean/SixBirdsFoundationsV/Laws/E8ControlPrice.lean}:153 \\
\Cref{lem:e8} & E8, E8.1, consistency lemma \newline catalog line 253 & \path{E8_ControlPrice} (theorem) \newline \path{lean/SixBirdsFoundationsV/Laws/E8ControlPrice.lean}:1608 \newline \path{E8_ComponentShadowPrice} (theorem) \newline \path{lean/SixBirdsFoundationsV/Laws/E8ControlPrice.lean}:2007 \newline \path{E8_CompressedSummaryLawfulness} (theorem) \newline \path{lean/SixBirdsFoundationsV/Laws/E8ControlPrice.lean}:2037 \newline \path{E8_1_AuditCapture} (theorem) \newline \path{lean/SixBirdsFoundationsV/Laws/E8ControlPrice.lean}:2050 \newline \path{E8_1_CaptureFalsifier} (theorem) \newline \path{lean/SixBirdsFoundationsV/Laws/E8ControlPrice.lean}:2341 \newline \path{E8_1_MetaAuditBoundedByE2} (theorem) \newline \path{lean/SixBirdsFoundationsV/Laws/E8ControlPrice.lean}:2365 \newline \path{BudgetInflationAudit.lineageMustMatchClassifiedMove} (theorem) \newline \path{lean/SixBirdsFoundationsV/Laws/E8ControlPrice.lean}:637 \\
\Cref{tmpl:e10} & E10--E10.2, conditional classification \newline catalog line 293 & \path{E10_CognitiveDemarcation} (theorem) \newline \path{lean/SixBirdsFoundationsV/Laws/E10CognitiveDemarcation.lean}:1771 \newline \path{E10_ScheduleTrapNull} (theorem) \newline \path{lean/SixBirdsFoundationsV/Laws/E10CognitiveDemarcation.lean}:2027 \newline \path{E10_ScheduleTrapExcludesLowerPriority} (theorem) \newline \path{lean/SixBirdsFoundationsV/Laws/E10CognitiveDemarcation.lean}:1963 \newline \path{E10_DecodableCorrelate} (theorem) \newline \path{lean/SixBirdsFoundationsV/Laws/E10CognitiveDemarcation.lean}:2012 \newline \path{E10_1_Intention} (theorem) \newline \path{lean/SixBirdsFoundationsV/Laws/E10CognitiveDemarcation.lean}:2042 \newline \path{E10_1_PostHocIntentionFalsifier} (theorem) \newline \path{lean/SixBirdsFoundationsV/Laws/E10CognitiveDemarcation.lean}:2186 \newline \path{E10_1_PostHocIntentionExcludesLowerPriority} (theorem) \newline \path{lean/SixBirdsFoundationsV/Laws/E10CognitiveDemarcation.lean}:2155 \newline \path{E10_2_Goal} (theorem) \newline \path{lean/SixBirdsFoundationsV/Laws/E10CognitiveDemarcation.lean}:2201 \newline \path{E10_2_RewardProxyNotGoal} (theorem) \newline \path{lean/SixBirdsFoundationsV/Laws/E10CognitiveDemarcation.lean}:2314 \newline \path{E10_2_RewardProxyOnlyExcludesLowerPriority} (theorem) \newline \path{lean/SixBirdsFoundationsV/Laws/E10CognitiveDemarcation.lean}:2291 \\
\Cref{tmpl:e11} & E11, conditional classification \newline catalog line 323 & \path{E11_StackActivity} (theorem) \newline \path{lean/SixBirdsFoundationsV/Laws/E11InstitutionalRewrite.lean}:811 \newline \path{E11_Constitutive} (theorem) \newline \path{lean/SixBirdsFoundationsV/Laws/E11InstitutionalRewrite.lean}:853 \newline \path{E11_ConditioningOnly} (theorem) \newline \path{lean/SixBirdsFoundationsV/Laws/E11InstitutionalRewrite.lean}:900 \newline \path{E11_WashoutInert} (theorem) \newline \path{lean/SixBirdsFoundationsV/Laws/E11InstitutionalRewrite.lean}:941 \newline \path{E11_NCTDObstruction} (theorem) \newline \path{lean/SixBirdsFoundationsV/Laws/E11InstitutionalRewrite.lean}:966 \newline \path{E11_StatusPartition} (theorem) \newline \path{lean/SixBirdsFoundationsV/Laws/E11InstitutionalRewrite.lean}:973 \\
\Cref{tmpl:e12} & E12, template \newline catalog line 343 & \path{E12_Individuation} (theorem) \newline \path{lean/SixBirdsFoundationsV/Laws/E12Individuation.lean}:711 \newline \path{E12_StatusPartition} (theorem) \newline \path{lean/SixBirdsFoundationsV/Laws/E12Individuation.lean}:787 \\
\Cref{prop:e12-singleton} & E12, template, singleton case \newline catalog line 343 & \path{Instances.singletonCandidates} (def) \newline \path{lean/SixBirdsFoundationsV/Instances/E12Singleton.lean}:5 \newline \path{Instances.E12_singleton_maximal_from_tests} (theorem) \newline \path{lean/SixBirdsFoundationsV/Instances/E12Singleton.lean}:12 \newline \path{Instances.spikeSelfMaintained} (theorem) \newline \path{lean/SixBirdsFoundationsV/Instances/SpikeModel.lean}:192 \\
\Cref{tmpl:e13} & E13--E13.4, template \newline catalog line 367 & \path{E13_Communication} (theorem) \newline \path{lean/SixBirdsFoundationsV/Laws/E13RepairTransport.lean}:1760 \newline \path{E13_TeachingCapacity} (theorem) \newline \path{lean/SixBirdsFoundationsV/Laws/E13RepairTransport.lean}:1773 \newline \path{E13_CoercionNull} (theorem) \newline \path{lean/SixBirdsFoundationsV/Laws/E13RepairTransport.lean}:1786 \newline \path{E13_SymbolicRepair} (theorem) \newline \path{lean/SixBirdsFoundationsV/Laws/E13RepairTransport.lean}:1799 \newline \path{E13_StatusPartition} (theorem) \newline \path{lean/SixBirdsFoundationsV/Laws/E13RepairTransport.lean}:1812 \newline \path{E13_CoercionNullExcludesLowerPriority} (theorem) \newline \path{lean/SixBirdsFoundationsV/Laws/E13RepairTransport.lean}:1992 \newline \path{Instances.transportContextZero} (def) \newline \path{lean/SixBirdsFoundationsV/Instances/E13Contexts.lean}:5 \newline \path{Instances.transportContextOne} (def) \newline \path{lean/SixBirdsFoundationsV/Instances/E13Contexts.lean}:10 \newline \path{Instances.E13_two_contexts_distinct} (theorem) \newline \path{lean/SixBirdsFoundationsV/Instances/E13Contexts.lean}:18 \newline \path{Instances.E13_two_contexts_members} (theorem) \newline \path{lean/SixBirdsFoundationsV/Instances/E13Contexts.lean}:22 \\
\Cref{tmpl:e14} & E14, conditional classification \newline catalog line 392 & \path{E14_ReconsolidationStatus} (theorem) \newline \path{lean/SixBirdsFoundationsV/Laws/E14Reconsolidation.lean}:1767 \newline \path{E14_LawfulConflictTrichotomy} (theorem) \newline \path{lean/SixBirdsFoundationsV/Laws/E14Reconsolidation.lean}:2051 \newline \path{E14_OutcomeCollisionExcludesLawful} (theorem) \newline \path{lean/SixBirdsFoundationsV/Laws/E14Reconsolidation.lean}:2105 \newline \path{E14_RecordRepair} (theorem) \newline \path{lean/SixBirdsFoundationsV/Laws/E14Reconsolidation.lean}:2217 \newline \path{E14_RecordCoarsening} (theorem) \newline \path{lean/SixBirdsFoundationsV/Laws/E14Reconsolidation.lean}:2238 \newline \path{E14_StatusedUnresolvedChargedToLambda} (theorem) \newline \path{lean/SixBirdsFoundationsV/Laws/E14Reconsolidation.lean}:2259 \newline \path{E14_SilentRewriteFalsifier} (theorem) \newline \path{lean/SixBirdsFoundationsV/Laws/E14Reconsolidation.lean}:2280 \newline \path{E14_RetrievalWithoutTransportControl} (theorem) \newline \path{lean/SixBirdsFoundationsV/Laws/E14Reconsolidation.lean}:2322 \newline \path{E14_LaunderedProvenanceDefect} (theorem) \newline \path{lean/SixBirdsFoundationsV/Laws/E14Reconsolidation.lean}:2365 \\
\Cref{tmpl:e15} & E15, conditional classification \newline catalog line 412 & \path{E15_OfflineReclosureStatus} (theorem) \newline \path{lean/SixBirdsFoundationsV/Laws/E15OfflineReclosure.lean}:2043 \newline \path{E15_DeficitAlternation} (theorem) \newline \path{lean/SixBirdsFoundationsV/Laws/E15OfflineReclosure.lean}:2135 \newline \path{E15_OnlineSufficientNoOfflineRequired} (theorem) \newline \path{lean/SixBirdsFoundationsV/Laws/E15OfflineReclosure.lean}:2182 \newline \path{E15_OfflinePermittedOutsideDeficit} (theorem) \newline \path{lean/SixBirdsFoundationsV/Laws/E15OfflineReclosure.lean}:2216 \newline \path{E15_SkippedOfflineCascade} (theorem) \newline \path{lean/SixBirdsFoundationsV/Laws/E15OfflineReclosure.lean}:2241 \newline \path{E15_DecorativeOfflineFalsifier} (theorem) \newline \path{lean/SixBirdsFoundationsV/Laws/E15OfflineReclosure.lean}:2266 \newline \path{E15_DutyCycleMismatchFalsifier} (theorem) \newline \path{lean/SixBirdsFoundationsV/Laws/E15OfflineReclosure.lean}:2300 \newline \path{E15_DeficitOnlineFalsifiesNecessity} (theorem) \newline \path{lean/SixBirdsFoundationsV/Laws/E15OfflineReclosure.lean}:2334 \newline \path{E15_DeficitCounterexampleExcludesLowerPriority} (theorem) \newline \path{lean/SixBirdsFoundationsV/Laws/E15OfflineReclosure.lean}:2369 \newline \path{E15_OffInventoryFlowCannotCertifyOffline} (theorem) \newline \path{lean/SixBirdsFoundationsV/Laws/E15OfflineReclosure.lean}:2456 \\
\Cref{tmpl:e16} & E16, conditional classification \newline catalog line 430 & \path{E16_AdaptabilityStatus} (theorem) \newline \path{lean/SixBirdsFoundationsV/Laws/E16Adaptability.lean}:2126 \newline \path{E16_CoherentAdaptability} (theorem) \newline \path{lean/SixBirdsFoundationsV/Laws/E16Adaptability.lean}:2235 \newline \path{E16_CurrentAuditBlindness} (theorem) \newline \path{lean/SixBirdsFoundationsV/Laws/E16Adaptability.lean}:2251 \newline \path{E16_RouteManufacturesRepairCapacity} (theorem) \newline \path{lean/SixBirdsFoundationsV/Laws/E16Adaptability.lean}:2272 \newline \path{E16_ArtifactNotAdaptability} (theorem) \newline \path{lean/SixBirdsFoundationsV/Laws/E16Adaptability.lean}:2288 \newline \path{E16_FlatNoAdaptability} (theorem) \newline \path{lean/SixBirdsFoundationsV/Laws/E16Adaptability.lean}:2302 \newline \path{E16_FlattenableNotAdaptability} (theorem) \newline \path{lean/SixBirdsFoundationsV/Laws/E16Adaptability.lean}:2316 \newline \path{E16_CurrentizableIsLegibleSlack} (theorem) \newline \path{lean/SixBirdsFoundationsV/Laws/E16Adaptability.lean}:2331 \newline \path{E16_DissipativeNotAdaptability} (theorem) \newline \path{lean/SixBirdsFoundationsV/Laws/E16Adaptability.lean}:2359 \newline \path{E16_BoundedPerturbationFailureRejected} (theorem) \newline \path{lean/SixBirdsFoundationsV/Laws/E16Adaptability.lean}:2374 \newline \path{E16_SupportConfoundExcludesLowerPriority} (theorem) \newline \path{lean/SixBirdsFoundationsV/Laws/E16Adaptability.lean}:2421 \newline \path{E16_LoopAsymmetryAnchor} (theorem) \newline \path{lean/SixBirdsFoundationsV/Laws/E16Adaptability.lean}:2481 \newline \path{E16_SwappedRouteTrialPopulationRejected} (theorem) \newline \path{lean/SixBirdsFoundationsV/Laws/E16Adaptability.lean}:457 \newline \path{E16_FlatSecondArmUsesClaimLinkedPair} (theorem) \newline \path{lean/SixBirdsFoundationsV/Laws/E16Adaptability.lean}:1379 \newline \path{E16_UnrelatedEqualPairCannotReplaceClaimDifference} (theorem) \newline \path{lean/SixBirdsFoundationsV/Laws/E16Adaptability.lean}:1406 \\

\end{longtable}
}

\Cref{tab:catalog-index} gives, for each definition and entry, its name in the theorem catalog and the
status under which the main paper states it. The catalog states every entry E1--E16 as a theorem
about arbitrary systems; that form is not a proof grade. Where the catalog and the main paper differ, the main paper's status is the one
supported by the Lean statements.

{\small
\begin{longtable}{@{}l>{\raggedright\arraybackslash}p{5.6cm}>{\raggedright\arraybackslash}p{5.6cm}r@{}}
\caption{Catalog name and status in the main paper.}\label{tab:catalog-index}\\
\toprule
\textbf{Entry} & \textbf{Catalog name} & \textbf{Main paper} & \textbf{Line} \\
\midrule
\endfirsthead
\toprule
\textbf{Entry} & \textbf{Catalog name} & \textbf{Main paper} & \textbf{Line} \\
\midrule
\endhead
\bottomrule
\endfoot
D1 & Repair join & definition; general theorem & 54 \\
D2 & Predictive surplus, quantified & definition; conditional general theorem & 56 \\
D3 & Carried record; carried instrument & definition; consistency lemma & 60 \\
D4 & Endogenous closure system, E-system & definition; consistency lemma & 62 \\
D5 & Closed-loop scope & definition; consistency lemma & 66 \\
D6 & Exposure cost; probe economy & definition; consistency lemma & 68 \\
E1 & Internalization Law & law, weak instance & 80 \\
E2 & Bounded Reflexivity Law & law, weak instance & 109 \\
E3 & Self-Maintaining Reclosure Law & law, weak instance & 133 \\
E4 & Repair Compilation Law & certificate specification & 157 \\
E5 & Reclosure Collapse Law & definition & 181 \\
E6 & Attention Law & general theorem & 205 \\
E7 & Alarm Law & conditional classification & 229 \\
E8 & Control-Price Law & definition & 253 \\
E9 & Curiosity Law & general theorem & 273 \\
E10 & Cognitive Demarcation Law & conditional classification & 293 \\
E11 & Institutional Rewrite Law & conditional classification & 323 \\
E12 & Individuation Law & law, weak instance & 343 \\
E13 & Repair Transport Law & law, weak instance & 367 \\
E14 & Reconsolidation Law & conditional classification & 392 \\
E15 & Offline Reclosure Law & conditional classification & 412 \\
E16 & Adaptability Law & conditional classification & 430 \\

\end{longtable}
}
